\section{Derivation of the shared-site identity}
\label{app:shared-site-derivation}

This appendix derives Eqs.~\eqref{eq:common-quadratic-overlap}--
\eqref{eq:common-source-with-residual}.  The assumptions and operator
expansions are those of Section~\ref{sec:shared-site-common-form}.  All
lower symbols below are evaluated with the same rank-one spin state
$\rho=(\Id_2+S^i\sigma^i)/2$ on the two physical sites, while the
auxiliary-space matrix order is kept fixed.

From Eq.~\eqref{eq:common-time-coefficients},
\begin{equation}
A^{(1)}=-2\ii(X_1-X_2)P.
\label{eq:app-A1-permutation}
\end{equation}
At coincident spins,
$P(\rho\otimes\rho)=(\rho\otimes\rho)P=\rho\otimes\rho$, so
$\langle A^{(1)}\rangle=0$.  Direct substitution in the two overlap
products gives
\begin{equation}
\mathcal D^{(+)}_2(\boldsymbol S,\boldsymbol S)
=\mathcal D^{(-)}_2(\boldsymbol S,\boldsymbol S)
=2\ii\left[(X^2)^\downarrow-U^2\right]
=2\ii\Xi.
\label{eq:app-D2-coincident}
\end{equation}
Thus the coincident order-$\epsilon^2$ contribution cancels in
$\widehat{\mathcal D}^{(+)}-\widehat{\mathcal D}^{(-)}$.

To obtain its first spatial variation, choose at one point a fixed physical
basis in which
\begin{equation}
\rho=\begin{pmatrix}1&0\\0&0\end{pmatrix},
\qquad
\partial_x\rho=\begin{pmatrix}0&p\\q&0\end{pmatrix},
\qquad
X=\sum_{r,s=0}^1X_{rs}\otimes|r\rangle\langle s|.
\label{eq:app-rho-block-basis}
\end{equation}
The basis is held fixed when differentiating.  The projector identity
$\rho\,\partial_x\rho+(\partial_x\rho)\rho=\partial_x\rho$ removes the
diagonal components of $\partial_x\rho$.  In this basis
\begin{equation}
U=X_{00},
\qquad
\Xi=X_{01}X_{10},
\label{eq:app-U-Xi-blocks}
\end{equation}
and differentiation with respect to the spin state gives
\begin{equation}
\partial_x\Xi
=p(X_{11}-X_{00})X_{10}
+qX_{01}(X_{11}-X_{00}).
\label{eq:app-Xi-tangent}
\end{equation}
Let $\delta_R$ and $\delta_L$ denote directional variations with respect to
the right and left physical density matrices, respectively.  Their action on
the quadratic overlap coefficients is
\begin{align}
\delta_R\mathcal D^{(+)}_2[\partial_x\rho]
&=2\ii p(X_{11}-X_{00})X_{10},
\nonumber\\
\delta_L\mathcal D^{(-)}_2[\partial_x\rho]
&=2\ii qX_{01}(X_{11}-X_{00}).
\label{eq:app-D2-tangent}
\end{align}
The minus sign in the definition of the overlap source is compensated by the
Taylor expansion of the left neighbour,
$\rho(x-\epsilon)=\rho(x)-\epsilon\partial_x\rho+\cdots$.  Therefore the two
terms in Eq.~\eqref{eq:app-D2-tangent} add and reproduce
Eq.~\eqref{eq:common-quadratic-Taylor}.

It remains to evaluate the coincident order-$\epsilon^3$ term.  Using
$\langle A^{(1)}\rangle=0$, one finds
\begin{align}
\mathcal D^{(+)}_3-\mathcal D^{(-)}_3
={}&\left\langle
A^{(2)}X_1-X_2A^{(2)}+A^{(1)}Y_1-Y_2A^{(1)}
\right\rangle
-[\langle A^{(2)}\rangle,U].
\label{eq:app-D3-start}
\end{align}
We separate the four terms in $A^{(2)}$ according to their origin.
The exchange part gives
\begin{equation}
(\mathcal D^{(+)}_3-\mathcal D^{(-)}_3)_{\rm ex}
=2\ii[\Xi,U]
-\ii\left\langle X_1[2P,X_2X_1]\right\rangle.
\label{eq:app-D3-exchange}
\end{equation}
The terms containing the explicit second spatial coefficient cancel only
after both its occurrence in $A^{(2)}$ and the products with $A^{(1)}$ are
retained:
\begin{align}
(\mathcal D^{(+)}_3-\mathcal D^{(-)}_3)_Y
&=-2\ii\left\langle
([Y_1,X_2]+[X_1,Y_2])P
\right\rangle
\nonumber\\
&=-2\ii\left([Y^\downarrow,U]+[U,Y^\downarrow]\right)=0.
\label{eq:app-D3-Y}
\end{align}
The spectral-derivative contribution is
\begin{equation}
(\mathcal D^{(+)}_3-\mathcal D^{(-)}_3)_Z
=-\ii\left\langle X_1(-Z_1+Z_2)\right\rangle.
\label{eq:app-D3-Z}
\end{equation}

For the Hamiltonian correction define
\begin{equation}
K_h:=\langle[h^{(1)},X_2]\rangle.
\label{eq:app-Kh}
\end{equation}
The triplet condition~\eqref{eq:common-triplet-condition} implies
\begin{align}
\left\langle
(X_1+X_2)[h^{(1)},X_1+X_2]
\right\rangle&=0,
\label{eq:app-h-collective}\\
\langle[h^{(1)},X_2X_1]\rangle&=[K_h,U].
\label{eq:app-h-product}
\end{align}
The first identity follows because $X_1+X_2$ commutes with the site
permutation and therefore preserves the triplet containing the product
state.  For the second identity, use the basis of
Eq.~\eqref{eq:app-rho-block-basis}.  If
$a_h=(h^{(1)})_{00,01}$ and $b_h=(h^{(1)})_{01,00}$, the triplet condition
gives
\begin{equation}
K_h=a_hX_{10}-b_hX_{01},
\qquad
\langle[h^{(1)},X_2X_1]\rangle
=a_h[X_{10},X_{00}]-b_h[X_{01},X_{00}],
\label{eq:app-h-blocks}
\end{equation}
which is Eq.~\eqref{eq:app-h-product}.  The complete Hamiltonian contribution
then reduces to
\begin{equation}
(\mathcal D^{(+)}_3-\mathcal D^{(-)}_3)_h
=-\ii\left\langle
X_1[h^{(1)},X_1+X_2]
\right\rangle.
\label{eq:app-D3-h}
\end{equation}

Adding Eqs.~\eqref{eq:app-D3-exchange}, \eqref{eq:app-D3-Y},
\eqref{eq:app-D3-Z}, and \eqref{eq:app-D3-h}, and using the definition
\eqref{eq:common-Sutherland-residual}, gives
\begin{equation}
\mathcal D^{(+)}_3(\boldsymbol S,\boldsymbol S)
-\mathcal D^{(-)}_3(\boldsymbol S,\boldsymbol S)
=2\ii[\Xi,U]-\ii\langle X_1\mathcal B\rangle.
\label{eq:app-D3-final}
\end{equation}
Together with the Taylor term~\eqref{eq:common-quadratic-Taylor}, this yields
\begin{equation}
\mathcal C
=2\ii\left(\partial_x\Xi+[\Xi,U]\right)
-\ii\langle X_1\mathcal B\rangle,
\label{eq:app-source-final}
\end{equation}
which is Eq.~\eqref{eq:common-source-with-residual}.  No
Cayley--Hamilton relation in the auxiliary space has been used, so the
argument allows arbitrary finite auxiliary dimension.  The two-dimensional
physical spin space enters through the rank-one spin-$\tfrac12$ projector and
through Eq.~\eqref{eq:app-h-blocks}.
